\documentclass[11pt]{article}
\usepackage[margin=25mm]{geometry}
\usepackage{fontspec}
\setmainfont{Latin Modern Roman}
\usepackage{amsmath,amssymb,mathtools}
\usepackage{unicode-math}
\setmathfont{Latin Modern Math}
\usepackage{xurl,hyperref}
\hypersetup{colorlinks=true,urlcolor=blue}
\setcounter{secnumdepth}{0}
\setlength{\parindent}{0pt}
\setlength{\parskip}{5pt}
\setlength{\emergencystretch}{3em}
\newcommand{\sref}[2]{\hyperlink{source:#1:#2}{[S#2]}}
\newcommand{\cataloguescope}[1]{\paragraph{Recorded target.}#1}
\newcommand{\R}{\mathbb R}
\newcommand{\C}{\mathbb C}
\newcommand{\N}{\mathbb N}
\newcommand{\Z}{\mathbb Z}
\newcommand{\Q}{\mathbb Q}
\newcommand{\F}{\mathbb F}
\newcommand{\eps}{\varepsilon}
\begin{document}
\section{Four-Dimensional Standard-Subspace Nets with Minimal Length}\label{statement}

\subsection{Definitions and mathematical statement}

\paragraph{Shared notation.}$\mathbb N=\{1,2,\ldots\}$; the standard number fields have their usual meanings. All additional conventions are stated locally.


\paragraph{Definitions and conventions.}
A closed real subspace $H$ of a complex Hilbert space is standard when $H\cap iH=\{0\}$ and $\overline{H+iH}$ is the whole space. A local net assigns such subspaces to spacetime regions, is isotone under inclusion, and makes spacelike-separated subspaces orthogonal for the symplectic form $\operatorname{Im}\langle\cdot,\cdot\rangle$. A double cone is the intersection of a future and a past light cone in Minkowski spacetime.

\paragraph{Target retained from the canonical source catalogue.}
\paragraph{Catalogue wording retained for traceability.}
Does there exist a natural local net of standard subspaces on double cones in Minkowski spacetime $\mathbb R^4$ with positive minimal length?

\paragraph{Source scope and qualification.}
The word natural and the precise minimal-length condition are not specified in the catalogue extraction. In particular, no particular double-cone radius threshold is invented as the original definition.
\subsection{Short English statement}
Can a natural local quantum-subspace net in four-dimensional Minkowski space have a positive minimum localization length?
\subsection{Research attempt}
\paragraph{Model, date and investigated scope.}
OpenAI GPT-6 Astra (configured model \texttt{gpt-6-astra}), ultra reasoning; generated 11 October 2026. Tests a naive radius-threshold net and proves it loses standardness and can lose additivity.

\paragraph{Sources and literature check.}
The complete supplied definition, including source qualifications and the dated literature assessment, was reviewed. The calculation below is an elementary restricted attempt motivated by [S2]; it does not claim novelty or a fresh literature-status determination.

\paragraph{Context and investigated scope.} The complete supplied definition, including source qualifications and the dated literature assessment, was reviewed. The calculation below is an elementary restricted attempt motivated by [S2]; it does not claim novelty or a fresh literature-status determination.

\paragraph{A tempting threshold construction and its defect.} Start with a local isotone real-subspace net $K(O)$ and choose a radius $\ell>0$ for a restricted family of nested standard double cones. Define $H(O)=0$ for radius below $\ell$ and $H(O)=K(O)$ above it. On that nested family isotony and locality survive: replacing subspaces by zero cannot violate symplectic orthogonality, and crossing the threshold only enlarges a subspace.

\paragraph{Standardness immediately fails below threshold.} In a nonzero ambient complex Hilbert space, the zero real subspace satisfies $0\cap i0=0$ but $\overline{0+i0}=0\ne\mathcal H$, so it is not standard. Thus this artificial cutoff does not meet a literal requirement of standard subspaces for every double cone. It can also break additivity: if a large region is covered by small regions assigned zero, their generated subspace is zero while the large-region assignment may be nonzero.

\paragraph{What remains to formulate and prove.} The source's positive minimal length may refer to a subtler localization property, and the catalogue does not define it precisely. The preceding construction is therefore a failed toy route, not a contradiction in the original problem. One must recover the intended minimal-length condition and covariance requirements and then construct a four-dimensional net satisfying them without silently weakening standardness or additivity. No naturalness criterion or Poincaré-covariant extension of the radius cutoff is claimed.

\paragraph{Verification and adversarial audit.}
Checked zero-subspace standardness, nested isotony/locality and cover-additivity failure; did not impose the toy radius threshold as the source definition. Model self-audit only; no independent human proof checking. Prompts A and B were applied to this canonical record and its qualified source scope. Checked zero-subspace standardness, nested isotony/locality and cover-additivity failure; did not impose the toy radius threshold as the source definition. Retained unresolved status and zero original-target progress; The source minimal-length convention and a natural covariant standard-subspace construction meeting it are not established.

Coordinator checked failure of standardness for the zero real subspace and the stated cutoff additivity obstruction. This additional check was model-assisted, not independent human verification.

\paragraph{Tools, process and contributions.}
Python 3 was used for file preparation and structural checks; no numerical experiment or theorem prover was used. Contribution prompt A was followed by prompt B. GPT-6 Astra with ultra reasoning drafted the research and performed its stated self-audit; any additional coordinator contribution is identified in the verification and audit. A GPT-6-based Codex coordinator assembled the manuscript and checked repository integration; its exact variant and reasoning configuration were not exposed. No human mathematical contribution, independent human proof check or formal proof certification is supplied. Complete derivations and the consulted source references are included here; private reasoning transcripts are not published.

\paragraph{Final status and first remaining gap.}
\textbf{UNRESOLVED}. The source minimal-length convention and a natural covariant standard-subspace construction meeting it are not established.

\section{Completion Estimate}
\noindent\textbf{COMPLETION: 0\%}
The calculation is an auxiliary or restricted attempt and establishes no newly verified portion of the full source target. The source minimal-length convention and a natural covariant standard-subspace construction meeting it are not established.

\subsection{Sources}
\begin{itemize}
\item[\hypertarget{source:statement:1}{S1}] ULAM.AI and the UnsolvedMath Contributors, UnsolvedMath, catalogue entry. Catalogue curation under CC BY 4.0; definitions and English exposition added locally.
\url{https://huggingface.co/datasets/ulamai/UnsolvedMath}.
\item[\hypertarget{source:statement:2}{S2}] Subfactors and Conformal Field Theory (2015), Oberwolfach Reports, DOI 10.4171/owr/2015/16, pp. 849--926. Primary report accessed 2026-10-11; relevant PDF pages 8, 9.
\url{https://doi.org/10.4171/owr/2015/16}.
\item[\hypertarget{source:statement:3}{S3}] Lechner--Longo, localization in standard spaces (catalogue-listed context; not independently consulted in this preparation).
\url{https://arxiv.org/abs/1403.1226}.
\item[\hypertarget{source:statement:4}{S4}] published localization results (catalogue-listed context; not independently consulted in this preparation).
\url{https://doi.org/10.1007/s00220-014-2199-2}.
\end{itemize}
\label{end:statement}
\end{document}
